\subsection{Vanishing of the kernel in first residue characteristic zero}\label{sec:second-three-charzero}
When the first residue field has characteristic zero, the results of Sections~\ref{sec:second-two-tame} and~\ref{sec:second-two-wild} also give the vanishing of the kernel of the reduced norm.

\begin{corollary}\label{cor:second-three-charzero}
For every three-dimensional local field $F$ whose first residue field has characteristic zero, and every central division algebra $D$ over $F$, one has $\SK_2(D)=0$.
This includes both characteristic sequences $(0,0,0)$ and $(0,0,p)$.
\end{corollary}
\begin{proof}
By Cohen's structure theorem, write $F=k((t))$, where $k$ is the first residue field of $F$.
Let $B$ be the first residue division algebra of $D$, and put $L=Z(B)$.
The center $L$ is a two-dimensional local field, so Corollary~\ref{cor:review-two-local-sk1-zero} gives $\SK_1(B)=0$, which is the hypothesis of Proposition~\ref{prop:second-residue-zero}.
By that proposition, $L/k$ is cyclic and
\[
 \SK_2(D)\simeq\ker\bigl[K_2(B)_{\Gal(L/k)}
  \xrightarrow{\ N_{L/k}\Nrd_{B/L,2}\ }K_2(k)\bigr].
\]
The equivariant map $\Nrd_{B/L,2}\colon K_2(B)\to K_2(L)$ is an isomorphism by Theorem~\ref{thm:second-two-vanishing}, and cyclic Hilbert 90 for $K_2$ \cite[Theorem 14.1]{MerkurjevSuslin1982} makes $N_{L/k}$ injective on $K_2(L)_{\Gal(L/k)}$.
Thus the kernel is zero.
\end{proof}

The following observation on the unramified case $D=B((t))$, used in Remark~\ref{rem:paper-sk2-dimension-four}, holds in every characteristic.

\begin{proposition}\label{prop:paper-laurent-sk2}
Let $B$ be a central division algebra over a field $k$, and put $F=k((t))$ and $D=B((t))$.
Then the map $\SK_1(B)\to\SK_2(D)$, $x\mapsto x\cdot\{t\}$, is well defined and injective.
\end{proposition}
\begin{proof}
Multiplication by $\{t\}$ commutes with reduced norms in every characteristic: $K_1(B)$ is generated by splitting transfers, and for $x=\theta_E(e)$ the projection formula and $\Nrd_2\theta_{E((t))}=N_{E((t))/F}$ give $\Nrd_{D,2}(x\cdot\{t\})=N_{E/k}(e)\cdot\{t\}=\Nrd_{B/k,1}(x)\cdot\{t\}$.
Hence multiplication by $\{t\}$ maps $\SK_1(B)$ into $\SK_2(D)$, and it is injective since the localization boundary $K_2(D)\to K_1(B)$ sends $x\cdot\{t\}$ to $\pm x$.
The same computation applies with $E(t)$ in place of $E((t))$.
\end{proof}
